\section{家庭、劳作与知识的早明证据}

本组材料把视线从朝廷决策移向家庭、劳作、环境与知识如何在文书中留下痕迹。土地、户籍、税役、河工、学校、宗教空间与地方灾情，往往只在征解、工程、诉讼、捐输或褒奖的片段中出现。材料的缺口并非无关紧要：它说明哪些人的行动更容易进入档案，哪些人的经验只能从地方记录和物质遗存中间接追索。

\subsection{户与田的记录方式}

\paragraph{1369（洪武二年）：中都（凤阳）开工建设} 这一锚点使家庭、劳动或地方资源进入可追查的历史时间。账本注明的把握范围是编年候选的年序可由官修与后出材料交叉；具体月日、参与者、战果或数字必须回查所列材料，不能将单一叙事当作定论。《太祖实录》洪武二年条；《明史·本纪》及相关列传；诏令、奏疏或会典版本 提供了定位事件的材料链，却分别反映编纂者、报送者和保存者的视角。其发生条件可概括为继承、官僚协调或皇权运作的压力使问题进入朝廷程序；随后可见的制度或社会影响是它影响后续人事与政策议程；动机和派系标签不能由单一叙事裁定。对于日常生活、性别与家庭分工、非精英劳作或未留名者，仍不能由这些材料直接补足。译写候选以编年材料定位；实录记载反映朝廷选择，崇祯期须另以奏疏、地方及敌对政权材料交叉。

\paragraph{1370（洪武三年）一月：阔克帖木儿围攻兰州但未能攻克} 这一锚点使家庭、劳动或地方资源进入可追查的历史时间。账本注明的把握范围是编年候选的年序可由官修与后出材料交叉；具体月日、参与者、战果或数字必须回查所列材料，不能将单一叙事当作定论。《太祖实录》洪武三年条；《明史·兵志》及相关列传；边镇、地方志或对方材料 提供了定位事件的材料链，却分别反映编纂者、报送者和保存者的视角。其发生条件可概括为前期边防、地方武装或跨区域资源调动的压力推动了该行动；随后可见的制度或社会影响是它改变了后续调兵、补给或地方秩序的条件；战果和伤亡须分开核对。对于日常生活、性别与家庭分工、非精英劳作或未留名者，仍不能由这些材料直接补足。译写候选以编年材料定位；实录记载反映朝廷选择，崇祯期须另以奏疏、地方及敌对政权材料交叉。

\paragraph{1370（洪武三年）：明军在巩昌击败了可克帖木儿，但未能俘获他} 这一锚点使家庭、劳动或地方资源进入可追查的历史时间。账本注明的把握范围是编年候选的年序可由官修与后出材料交叉；具体月日、参与者、战果或数字必须回查所列材料，不能将单一叙事当作定论。《太祖实录》洪武三年条；《明史·兵志》及相关列传；边镇、地方志或对方材料 提供了定位事件的材料链，却分别反映编纂者、报送者和保存者的视角。其发生条件可概括为前期边防、地方武装或跨区域资源调动的压力推动了该行动；随后可见的制度或社会影响是它改变了后续调兵、补给或地方秩序的条件；战果和伤亡须分开核对。对于日常生活、性别与家庭分工、非精英劳作或未留名者，仍不能由这些材料直接补足。译写候选以编年材料定位；实录记载反映朝廷选择，崇祯期须另以奏疏、地方及敌对政权材料交叉。

\paragraph{1370（洪武三年）六月5日：洪武皇帝授权明第一次科举考试} 这一锚点使家庭、劳动或地方资源进入可追查的历史时间。账本注明的把握范围是编年候选的年序可由官修与后出材料交叉；具体月日、参与者、战果或数字必须回查所列材料，不能将单一叙事当作定论。《太祖实录》洪武三年条；《明史·选举志》或《艺文志》；文集、碑刻或书目材料 提供了定位事件的材料链，却分别反映编纂者、报送者和保存者的视角。其发生条件可概括为制度、教育或知识传播的需要推动了相关行动或文本形成；随后可见的制度或社会影响是它提供制度和传播史的锚点，不能据此推断社会整体接受。对于日常生活、性别与家庭分工、非精英劳作或未留名者，仍不能由这些材料直接补足。译写候选以编年材料定位；实录记载反映朝廷选择，崇祯期须另以奏疏、地方及敌对政权材料交叉。

\paragraph{1370（洪武三年）六月10日：明军攻克颍昌} 这一锚点使家庭、劳动或地方资源进入可追查的历史时间。账本注明的把握范围是编年候选的年序可由官修与后出材料交叉；具体月日、参与者、战果或数字必须回查所列材料，不能将单一叙事当作定论。《太祖实录》洪武三年条；《明史·兵志》及相关列传；边镇、地方志或对方材料 提供了定位事件的材料链，却分别反映编纂者、报送者和保存者的视角。其发生条件可概括为前期边防、地方武装或跨区域资源调动的压力推动了该行动；随后可见的制度或社会影响是它改变了后续调兵、补给或地方秩序的条件；战果和伤亡须分开核对。对于日常生活、性别与家庭分工、非精英劳作或未留名者，仍不能由这些材料直接补足。译写候选以编年材料定位；实录记载反映朝廷选择，崇祯期须另以奏疏、地方及敌对政权材料交叉。

\paragraph{1370（洪武三年）六月：洪武帝取缔白莲教和摩尼教} 这一锚点使家庭、劳动或地方资源进入可追查的历史时间。账本注明的把握范围是编年候选的年序可由官修与后出材料交叉；具体月日、参与者、战果或数字必须回查所列材料，不能将单一叙事当作定论。《太祖实录》洪武三年条；《明史·本纪》及相关列传；诏令、奏疏或会典版本 提供了定位事件的材料链，却分别反映编纂者、报送者和保存者的视角。其发生条件可概括为继承、官僚协调或皇权运作的压力使问题进入朝廷程序；随后可见的制度或社会影响是它影响后续人事与政策议程；动机和派系标签不能由单一叙事裁定。对于日常生活、性别与家庭分工、非精英劳作或未留名者，仍不能由这些材料直接补足。译写候选以编年材料定位；实录记载反映朝廷选择，崇祯期须另以奏疏、地方及敌对政权材料交叉。

\subsection{军屯、工役与移动}

\paragraph{1370（洪武三年）：明代使用火药来增强地雷的爆炸力} 这一锚点使家庭、劳动或地方资源进入可追查的历史时间。账本注明的把握范围是编年候选的年序可由官修与后出材料交叉；具体月日、参与者、战果或数字必须回查所列材料，不能将单一叙事当作定论。《太祖实录》洪武三年条；《明会典》相关条目；地方志、黄册/契约/河工材料 提供了定位事件的材料链，却分别反映编纂者、报送者和保存者的视角。其发生条件可概括为财政征解、交通工程或货币支付的压力使相关措施进入政务；随后可见的制度或社会影响是其法定安排不等于地方执行，后续须以地域材料追踪负担与成效。对于日常生活、性别与家庭分工、非精英劳作或未留名者，仍不能由这些材料直接补足。译写候选以编年材料定位；实录记载反映朝廷选择，崇祯期须另以奏疏、地方及敌对政权材料交叉。

\paragraph{1370（洪武三年）：明代大炮的弹丸由石弹过渡为铁弹} 这一锚点使家庭、劳动或地方资源进入可追查的历史时间。账本注明的把握范围是编年候选的年序可由官修与后出材料交叉；具体月日、参与者、战果或数字必须回查所列材料，不能将单一叙事当作定论。《太祖实录》洪武三年条；《明史·选举志》或《艺文志》；文集、碑刻或书目材料 提供了定位事件的材料链，却分别反映编纂者、报送者和保存者的视角。其发生条件可概括为制度、教育或知识传播的需要推动了相关行动或文本形成；随后可见的制度或社会影响是它提供制度和传播史的锚点，不能据此推断社会整体接受。对于日常生活、性别与家庭分工、非精英劳作或未留名者，仍不能由这些材料直接补足。译写候选以编年材料定位；实录记载反映朝廷选择，崇祯期须另以奏疏、地方及敌对政权材料交叉。

\paragraph{1371（洪武四年）五月18日：明军攻占温州} 这一锚点使家庭、劳动或地方资源进入可追查的历史时间。账本注明的把握范围是编年候选的年序可由官修与后出材料交叉；具体月日、参与者、战果或数字必须回查所列材料，不能将单一叙事当作定论。《太祖实录》洪武四年条；《明史·兵志》及相关列传；边镇、地方志或对方材料 提供了定位事件的材料链，却分别反映编纂者、报送者和保存者的视角。其发生条件可概括为前期边防、地方武装或跨区域资源调动的压力推动了该行动；随后可见的制度或社会影响是它改变了后续调兵、补给或地方秩序的条件；战果和伤亡须分开核对。对于日常生活、性别与家庭分工、非精英劳作或未留名者，仍不能由这些材料直接补足。译写候选以编年材料定位；实录记载反映朝廷选择，崇祯期须另以奏疏、地方及敌对政权材料交叉。

\paragraph{1371（洪武四年）七月：明军攻克杭州} 这一锚点使家庭、劳动或地方资源进入可追查的历史时间。账本注明的把握范围是编年候选的年序可由官修与后出材料交叉；具体月日、参与者、战果或数字必须回查所列材料，不能将单一叙事当作定论。《太祖实录》洪武四年条；《明史·兵志》及相关列传；边镇、地方志或对方材料 提供了定位事件的材料链，却分别反映编纂者、报送者和保存者的视角。其发生条件可概括为前期边防、地方武装或跨区域资源调动的压力推动了该行动；随后可见的制度或社会影响是它改变了后续调兵、补给或地方秩序的条件；战果和伤亡须分开核对。对于日常生活、性别与家庭分工、非精英劳作或未留名者，仍不能由这些材料直接补足。译写候选以编年材料定位；实录记载反映朝廷选择，崇祯期须另以奏疏、地方及敌对政权材料交叉。

\paragraph{1371（洪武四年）八月3日：明盛将四川投降明朝} 这一锚点使家庭、劳动或地方资源进入可追查的历史时间。账本注明的把握范围是编年候选的年序可由官修与后出材料交叉；具体月日、参与者、战果或数字必须回查所列材料，不能将单一叙事当作定论。《太祖实录》洪武四年条；《明史·兵志》及相关列传；边镇、地方志或对方材料 提供了定位事件的材料链，却分别反映编纂者、报送者和保存者的视角。其发生条件可概括为前期边防、地方武装或跨区域资源调动的压力推动了该行动；随后可见的制度或社会影响是它改变了后续调兵、补给或地方秩序的条件；战果和伤亡须分开核对。对于日常生活、性别与家庭分工、非精英劳作或未留名者，仍不能由这些材料直接补足。译写候选以编年材料定位；实录记载反映朝廷选择，崇祯期须另以奏疏、地方及敌对政权材料交叉。

\paragraph{1371（洪武四年）：国子监注册学生达3,728人} 这一锚点使家庭、劳动或地方资源进入可追查的历史时间。账本注明的把握范围是编年候选的年序可由官修与后出材料交叉；具体月日、参与者、战果或数字必须回查所列材料，不能将单一叙事当作定论。《太祖实录》洪武四年条；《明史·本纪》及相关列传；诏令、奏疏或会典版本 提供了定位事件的材料链，却分别反映编纂者、报送者和保存者的视角。其发生条件可概括为继承、官僚协调或皇权运作的压力使问题进入朝廷程序；随后可见的制度或社会影响是它影响后续人事与政策议程；动机和派系标签不能由单一叙事裁定。对于日常生活、性别与家庭分工、非精英劳作或未留名者，仍不能由这些材料直接补足。译写候选以编年材料定位；实录记载反映朝廷选择，崇祯期须另以奏疏、地方及敌对政权材料交叉。

\subsection{环境、交通与地方物质}

\paragraph{1372（洪武五年）四月：明军在土拉河击败可克帖木儿} 这一锚点使家庭、劳动或地方资源进入可追查的历史时间。账本注明的把握范围是编年候选的年序可由官修与后出材料交叉；具体月日、参与者、战果或数字必须回查所列材料，不能将单一叙事当作定论。《太祖实录》洪武五年条；《明会典》相关条目；地方志、黄册/契约/河工材料 提供了定位事件的材料链，却分别反映编纂者、报送者和保存者的视角。其发生条件可概括为财政征解、交通工程或货币支付的压力使相关措施进入政务；随后可见的制度或社会影响是其法定安排不等于地方执行，后续须以地域材料追踪负担与成效。对于日常生活、性别与家庭分工、非精英劳作或未留名者，仍不能由这些材料直接补足。译写候选以编年材料定位；实录记载反映朝廷选择，崇祯期须另以奏疏、地方及敌对政权材料交叉。

\paragraph{1372（洪武五年）：明军在喀喇昆仑被击溃} 这一锚点使家庭、劳动或地方资源进入可追查的历史时间。账本注明的把握范围是编年候选的年序可由官修与后出材料交叉；具体月日、参与者、战果或数字必须回查所列材料，不能将单一叙事当作定论。《太祖实录》洪武五年条；《明史·兵志》及相关列传；边镇、地方志或对方材料 提供了定位事件的材料链，却分别反映编纂者、报送者和保存者的视角。其发生条件可概括为前期边防、地方武装或跨区域资源调动的压力推动了该行动；随后可见的制度或社会影响是它改变了后续调兵、补给或地方秩序的条件；战果和伤亡须分开核对。对于日常生活、性别与家庭分工、非精英劳作或未留名者，仍不能由这些材料直接补足。译写候选以编年材料定位；实录记载反映朝廷选择，崇祯期须另以奏疏、地方及敌对政权材料交叉。

\paragraph{1372（洪武五年）：明军攻克永昌，攻克居延} 这一锚点使家庭、劳动或地方资源进入可追查的历史时间。账本注明的把握范围是编年候选的年序可由官修与后出材料交叉；具体月日、参与者、战果或数字必须回查所列材料，不能将单一叙事当作定论。《太祖实录》洪武五年条；《明史·兵志》及相关列传；边镇、地方志或对方材料 提供了定位事件的材料链，却分别反映编纂者、报送者和保存者的视角。其发生条件可概括为前期边防、地方武装或跨区域资源调动的压力推动了该行动；随后可见的制度或社会影响是它改变了后续调兵、补给或地方秩序的条件；战果和伤亡须分开核对。对于日常生活、性别与家庭分工、非精英劳作或未留名者，仍不能由这些材料直接补足。译写候选以编年材料定位；实录记载反映朝廷选择，崇祯期须另以奏疏、地方及敌对政权材料交叉。

\paragraph{1372（洪武五年）：国子监注册学生突破1万人} 这一锚点使家庭、劳动或地方资源进入可追查的历史时间。账本注明的把握范围是编年候选的年序可由官修与后出材料交叉；具体月日、参与者、战果或数字必须回查所列材料，不能将单一叙事当作定论。《太祖实录》洪武五年条；《明史·本纪》及相关列传；诏令、奏疏或会典版本 提供了定位事件的材料链，却分别反映编纂者、报送者和保存者的视角。其发生条件可概括为继承、官僚协调或皇权运作的压力使问题进入朝廷程序；随后可见的制度或社会影响是它影响后续人事与政策议程；动机和派系标签不能由单一叙事裁定。对于日常生活、性别与家庭分工、非精英劳作或未留名者，仍不能由这些材料直接补足。译写候选以编年材料定位；实录记载反映朝廷选择，崇祯期须另以奏疏、地方及敌对政权材料交叉。

\paragraph{1372（洪武五年）：明代出现了专门为海军使用的大炮} 这一锚点使家庭、劳动或地方资源进入可追查的历史时间。账本注明的把握范围是编年候选的年序可由官修与后出材料交叉；具体月日、参与者、战果或数字必须回查所列材料，不能将单一叙事当作定论。《太祖实录》洪武五年条；《明史·选举志》或《艺文志》；文集、碑刻或书目材料 提供了定位事件的材料链，却分别反映编纂者、报送者和保存者的视角。其发生条件可概括为制度、教育或知识传播的需要推动了相关行动或文本形成；随后可见的制度或社会影响是它提供制度和传播史的锚点，不能据此推断社会整体接受。对于日常生活、性别与家庭分工、非精英劳作或未留名者，仍不能由这些材料直接补足。译写候选以编年材料定位；实录记载反映朝廷选择，崇祯期须另以奏疏、地方及敌对政权材料交叉。

\paragraph{1373（洪武六年）三月：洪武皇帝暂停科举考试} 这一锚点使家庭、劳动或地方资源进入可追查的历史时间。账本注明的把握范围是编年候选的年序可由官修与后出材料交叉；具体月日、参与者、战果或数字必须回查所列材料，不能将单一叙事当作定论。《太祖实录》洪武六年条；《明史·选举志》或《艺文志》；文集、碑刻或书目材料 提供了定位事件的材料链，却分别反映编纂者、报送者和保存者的视角。其发生条件可概括为制度、教育或知识传播的需要推动了相关行动或文本形成；随后可见的制度或社会影响是它提供制度和传播史的锚点，不能据此推断社会整体接受。对于日常生活、性别与家庭分工、非精英劳作或未留名者，仍不能由这些材料直接补足。译写候选以编年材料定位；实录记载反映朝廷选择，崇祯期须另以奏疏、地方及敌对政权材料交叉。

\subsection{学校、文本与知识秩序}

\paragraph{1373（洪武六年）十一月29日：明军在怀柔击败可克帖木儿} 这一锚点使家庭、劳动或地方资源进入可追查的历史时间。账本注明的把握范围是编年候选的年序可由官修与后出材料交叉；具体月日、参与者、战果或数字必须回查所列材料，不能将单一叙事当作定论。《太祖实录》洪武六年条；《明史·兵志》及相关列传；边镇、地方志或对方材料 提供了定位事件的材料链，却分别反映编纂者、报送者和保存者的视角。其发生条件可概括为前期边防、地方武装或跨区域资源调动的压力推动了该行动；随后可见的制度或社会影响是它改变了后续调兵、补给或地方秩序的条件；战果和伤亡须分开核对。对于日常生活、性别与家庭分工、非精英劳作或未留名者，仍不能由这些材料直接补足。译写候选以编年材料定位；实录记载反映朝廷选择，崇祯期须另以奏疏、地方及敌对政权材料交叉。

\paragraph{1373（洪武六年）：洪武皇帝将高丽的进贡限制为每三年一次} 这一锚点使家庭、劳动或地方资源进入可追查的历史时间。账本注明的把握范围是编年候选的年序可由官修与后出材料交叉；具体月日、参与者、战果或数字必须回查所列材料，不能将单一叙事当作定论。《太祖实录》洪武六年条；《明史·外国传》；《朝鲜王朝实录》、琉球《历代宝案》或海外档案 提供了定位事件的材料链，却分别反映编纂者、报送者和保存者的视角。其发生条件可概括为朝贡名义、边境安全与港口贸易的利益使交涉持续发生；随后可见的制度或社会影响是它为后续册封、互市或海防安排留下文书与跨政权比较线索。对于日常生活、性别与家庭分工、非精英劳作或未留名者，仍不能由这些材料直接补足。译写候选以编年材料定位；实录记载反映朝廷选择，崇祯期须另以奏疏、地方及敌对政权材料交叉。

\paragraph{1373（洪武六年）：明朝官员制定了中国历史上第一部《宅法》} 这一锚点使家庭、劳动或地方资源进入可追查的历史时间。账本注明的把握范围是编年候选的年序可由官修与后出材料交叉；具体月日、参与者、战果或数字必须回查所列材料，不能将单一叙事当作定论。《太祖实录》洪武六年条；《明史·本纪》及相关列传；诏令、奏疏或会典版本 提供了定位事件的材料链，却分别反映编纂者、报送者和保存者的视角。其发生条件可概括为继承、官僚协调或皇权运作的压力使问题进入朝廷程序；随后可见的制度或社会影响是它影响后续人事与政策议程；动机和派系标签不能由单一叙事裁定。对于日常生活、性别与家庭分工、非精英劳作或未留名者，仍不能由这些材料直接补足。译写候选以编年材料定位；实录记载反映朝廷选择，崇祯期须另以奏疏、地方及敌对政权材料交叉。

\paragraph{1374（洪武七年）：朝廷围绕卫所军户的编籍、屯田与调发整理本年军政文书，作为明初军事接收的可核查节点} 这一锚点使家庭、劳动或地方资源进入可追查的历史时间。账本注明的把握范围是来源导向的年度桥接条目：本年材料可定位相关政务线索；具体月日、发文机关、地域和执行结果仍须逐项回查，不能以制度文本或奏报替代实际效果。《太祖实录》洪武七年条；《明史·兵志》及相关列传；边镇、地方志或对方材料 提供了定位事件的材料链，却分别反映编纂者、报送者和保存者的视角。其发生条件可概括为前期边防、地方武装或跨区域资源调动的压力推动了该行动；随后可见的制度或社会影响是它改变了后续调兵、补给或地方秩序的条件；战果和伤亡须分开核对。对于日常生活、性别与家庭分工、非精英劳作或未留名者，仍不能由这些材料直接补足。桥接条目只标示可回查的年度问题与材料链；实录有中央选择性，崇祯期尤其需要奏疏、地方、朝鲜及满文材料交叉。

\paragraph{1374（洪武七年）：北元诸部、东北和西南的边报、招抚与防御命令持续进入本年政务，须按具体部众和地域复核} 这一锚点使家庭、劳动或地方资源进入可追查的历史时间。账本注明的把握范围是来源导向的年度桥接条目：本年材料可定位相关政务线索；具体月日、发文机关、地域和执行结果仍须逐项回查，不能以制度文本或奏报替代实际效果。《太祖实录》洪武七年条；《明史·外国传》；《朝鲜王朝实录》、琉球《历代宝案》或海外档案 提供了定位事件的材料链，却分别反映编纂者、报送者和保存者的视角。其发生条件可概括为朝贡名义、边境安全与港口贸易的利益使交涉持续发生；随后可见的制度或社会影响是它为后续册封、互市或海防安排留下文书与跨政权比较线索。对于日常生活、性别与家庭分工、非精英劳作或未留名者，仍不能由这些材料直接补足。桥接条目只标示可回查的年度问题与材料链；实录有中央选择性，崇祯期尤其需要奏疏、地方、朝鲜及满文材料交叉。

\paragraph{1374（洪武七年）：沿海朝贡、海禁和港口管理的文书构成本年海上联系线索，禁令不能直接推作私人航行停止} 这一锚点使家庭、劳动或地方资源进入可追查的历史时间。账本注明的把握范围是来源导向的年度桥接条目：本年材料可定位相关政务线索；具体月日、发文机关、地域和执行结果仍须逐项回查，不能以制度文本或奏报替代实际效果。《太祖实录》洪武七年条；《明史·外国传》；《朝鲜王朝实录》、琉球《历代宝案》或海外档案 提供了定位事件的材料链，却分别反映编纂者、报送者和保存者的视角。其发生条件可概括为朝贡名义、边境安全与港口贸易的利益使交涉持续发生；随后可见的制度或社会影响是它为后续册封、互市或海防安排留下文书与跨政权比较线索。对于日常生活、性别与家庭分工、非精英劳作或未留名者，仍不能由这些材料直接补足。桥接条目只标示可回查的年度问题与材料链；实录有中央选择性，崇祯期尤其需要奏疏、地方、朝鲜及满文材料交叉。
